\documentclass{zhvt-classic}

% 使用本地源流明体（GenRyuMin2 TW）按文件路径加载，不依赖系统字体名
% 版框（内框／界栏／书口侧栏／鱼尾）改用与版心印章同一印泥棕色：
%   印章图片 yshs.png 实色 = RGB(135,68,52) = HTML 874434
\usepackage{xcolor}
\usepackage{tikz}
% sealbrown 双注册：xcolor 供 TikZ/pgf，L3 供 cls 版框
\definecolor{sealbrown}{HTML}{874434}
\ExplSyntaxOn
\color_set:nnn { sealbrown } { HTML } { 874434 }
\ExplSyntaxOff

\zhvtset{
  grid color       = sealbrown,
  main font        = GenRyuMin2TW-SB,
  main font path   = C:/Users/Administrator/texmf/fonts/opentype/public/GenRyuMin2TW/,
jiazhu font      = GenRyuMin2TW-M,
  jiazhu font path = C:/Users/Administrator/texmf/fonts/opentype/public/GenRyuMin2TW/,
  % 页码框容纳 4 字：正常页「九十七」+前置空 1 格=4 行恰好；
  % 100+ 页「一百零一」共 5 行超出，末字「一/二/三」被框裁掉（2026-10-04 实测），故放宽到 5
  shukou page number chars = 5,
  % 书口书名框高从默认 6 字减到 4 字，把省出的 2 字高拨给 chapter 框，
  % 使 6 字书名+卷名放进 9 字高的 chapter 框后不再压线（2026-10-05）
  shukou title chars = 4,
}

\title{史记}[樗里子甘茂列传]
\author{汉·司马迁}
\maker{闲叟漫记·重製}


% ===== 索隐述赞样式（2026-11-08 四次定稿 · 愚大人亲手改定参数）=====
% ⚠ 竖排前提：自定义字体家族必须带 RawFeature={vertical:+vert}；
%   simhei.ttf 加 +vhal 会整体加载失败静默回退，只写 +vert；simfang 同理。
\newCJKfontfamily{\shuzanfont}[Path=C:/Windows/Fonts/,RawFeature={vertical:+vert}]{simhei.ttf}
\newCJKfontfamily{\shuzanbody}[Path=C:/Windows/Fonts/,RawFeature={vertical:+vert}]{simfang.ttf}
% 题头：152x30pt 圆角 6pt 棕底白字；baseline (0,2) 使整框落入界行格子不压版框线；
% 白字节点 x 不补偿、y +8.2pt（像素级墨迹实测 dx/dy≈0，L/R 留白约 20pt）。
\newcommand{\shuzanhead}[1]{%
  \begin{tikzpicture}[baseline={(0,2)}]
    \fill[sealbrown, rounded corners=6pt]
      (-76pt,-15pt) rectangle (76pt,15pt);
    \node[text=white, inner sep=0pt, anchor=center] at (0pt,8.2pt)
      {\shuzanfont\fontsize{23}{23}\selectfont #1};
  \end{tikzpicture}%
}
% ★ 修正书口卷名位置：chapter 框起始位置改回 cls 默认值，避免全局影响。
%   只在每个卷首页通过 \juan 命令里的空命令微调该页位置（安全，不会跨页污染）。
\ExplSyntaxOn
\AtBeginDocument{
  \dim_set:Nn \l__zhvt_shukou_chapter_title_start_dim
    { \l__zhvt_shukou_title_height_dim + \l__zhvt_font_size_dim * 2 }
  \dim_set:Nn \l__zhvt_shukou_page_number_start_dim
    { \l__zhvt_shukou_chapter_title_start_dim + \l__zhvt_shukou_chapter_title_height_dim }
}
\ExplSyntaxOff

\begin{document}

% ★★ 版心（书口）排布覆写：把书名放到鱼尾之下（2026-10-04）
%
%   cls 默认版心自上而下：书名（title 静态框，在鱼尾之上）→ 鱼尾 → 篇名（chapter 动态框）→ 页码
%   本模板改为：          鱼尾（位置不动）→ 书名 → 篇名 → 页码分割线 → 页码
%
%   手法：把 cls 放在 title 静态框里的书名，改放进 chapter 动态框（鱼尾之下）；
%         title 静态框留空。书名取 \l__zhvt_title_str，即 \title{...} 的第一参数，
%         所以换书时只改 \title，不必在此重复写书名。
%
%   ⚠ 三条必须遵守的写法（2026-10-04 实测）：
%     ① 汉字必须用 \hbox{...} 包住。flowfram 会把框内容装进一个宽 14.4mm（= 一个 line_sep
%        = 书口栏宽）的纵向 minipage 并按该宽度折行，而整页靠 PDF /Rotate 90 显示、
%        框内文字本身不变向 —— 不包 \hbox 时「围炉夜话」4 字会被逐字折成 4 行，
%        在版面上就成了横排 4 字并横向溢出版框。
%     ② 竖排下**一个 \hbox 就是一整列**。书名与篇名要同处一列，必须放进同一个 \hbox；
%        分成两个 \hbox 时，后一个会被当成新的一列、甩到隔壁栏去
%        （实测「序」跑到书口栏外侧 14.4mm 的正文列上）。中间的 \quad 在框内是横向空白，
%        显示出来正好是书名与篇名之间的纵向间隔。
%     ③ ★ 内容盒要写成 **与框等宽的定宽盒** `\hbox to \l__zhvt_shukou_chapter_title_height_dim{…\hfil}`。
%        实测：chapter 框是把内容盒沿纵向**居中**摆放的，内容宽度一变，书名就跟着漂：
%        自然宽度的 \hbox（书名 38.4mm）落在 x=157.2mm，加篇名后（57.6mm）落在 x=152.4mm，
%        目录页（无篇名）与正文页的书名就差半格。做成等宽盒后，无论有没有篇名，
%        内容固定从框顶起排，实测两页均落在 x=142.8mm，不再漂。
\ExplSyntaxOn
% 把书口 chapter 框（书名+篇名）起始位置下移 1 字高，避免 4 字卷名顶到鱼尾
% ⚠ cls 的 \__zhvt_calculate_dimensions: 在 \AtBeginDocument 里执行，会重算该 dim；
%   必须把覆写挂到 begin document 之后，否则会被二次覆盖（2026-10-05 实测）。
% (旧方案已移到导言区，直接修改 chapter 框起始位置)
% 卷标题命令 #1=书口篇名(鱼尾) #2=正文卷名（含卷序，卷名内空格用 ~）
%
% ★★ 目录页两项修订（2026-10-05，愚大人指出）★★
%   ① 页码起点统一对齐：
%      原状 —— 目录项写成 \par\hspace{4em}卷名\hspace{4em}页码，页码起点 = 卷名实宽 + 4em。
%      卷名实宽随字数变：「卷一~梁惠王上」6 字 +~、「卷七~离娄上」5 字 +~，
%      差 1 个字位 → 卷七~卷十 4 条的页码整体比其余 10 条高 27.21pt（一个字），看着上下不齐。
%      改法 —— 卷名补全角空格（宽 = 1 字 = 27.215pt）补足到 **7 字符位**再写进目录，
%      14 条页码起点遂全部落在同一条竖线上（起点对齐；末字随页码字数自然参差）。
%      ⚠ 只对目录里的副本 \l_tmpc_tl 补，正文卷名仍用 \l_tmpa_tl 原样 —— 正文像素不受影响。
%   ② 卷名整体上移半字位：前导 \hspace 由 4em 改 3.5em。
%      实测原状目录块在版框内偏下（天头留白 139.6pt / 地脚留白 110.7pt），
%      上移 0.5em = 13.6pt 后上下基本均分，视觉居中。
\NewDocumentCommand { \juan } { mm } {
	\clearpage
	\tl_set:Nn \l__zhvt_chapter_title_tl { #1 }
	\tl_set:Nn \l_tmpa_tl { #2 }
	\tl_set_eq:NN \l_tmpc_tl \l_tmpa_tl
	\int_step_inline:nn { 7 - \str_count:n { #2 } }
	  { \tl_put_right:Nn \l_tmpc_tl { 　 } }
	\tl_set:Nx \l_tmpb_tl { \zhnum { page } }
	\tl_put_left:Nn \l_tmpb_tl { \hspace { 4em } }
	\tl_put_left:Nn \l_tmpb_tl { \l_tmpc_tl }
	\tl_put_left:Nn \l_tmpb_tl { \par \hspace { 3.5em } }
	\addtocontents { toc } { \l_tmpb_tl }
	\skip_horizontal:n { \l__zhvt_font_size_dim * 4 }
	\l_tmpa_tl
	\par
	\thispagestyle { empty }
}
\cs_set_protected:Npn \__zhvt_set_shukou_content: {
	% ★ 鱼尾上方留白（title 静态框）嵌入闲叟钤印：
	%   整页 /Rotate 90，故此处把图用 \rotatebox{90} 转正为竖章；
	%   尺寸卡在版心条幅内（书内约 7mm 宽 × 33.8mm 高，留余量不压鱼尾）。
	\setstaticcontents * { title } {
		\vbox_to_ht:nn { \l__zhvt_baseline_skip_dim } {
			\vskip 14.2pt
			\hbox_to_wd:nn { \l__zhvt_shukou_title_height_dim } {
				\hfil
				\rotatebox [ origin = c ] { 90 } { \includegraphics [ width = 7mm , height = 33.8mm ] { yshs.png } }
				\hfil
			}
		}
	}
% ★ 版心（书口）统一改法（2026-11-16 愚大人定）：18pt＋框内居中＋书口文字同版框色
%   ① 原 27.2pt 时 7 字长篇名（如「廉颇蔺相如列传」）末字溢出 chapter 框、压住页码；
%      改 \fontsize{18pt}{18pt}，前后 \hfil 令字串在框内居中（不随篇名字数漂移）。
%   ② \vbox{\vskip 8.2pt} 定量下移，抵消缩字号后栏宽方向 8.2pt 上偏（像素实测 y 中心 543.3 = 原版）。
%      ⚠ \raisebox / \vspace* 在 flowfram 框内无效；换字号须同步换补偿：16pt→10.1、19pt→7.4。
%   ③ chapter 与 page~number 均置 \color{sealbrown}（HTML 874434），与版框／界行同色；正文仍黑。
%   ⚠ 页码字位、钤印、鱼尾、界行、正文一律未动；五部合订本同此口径。
	\setdynamiccontents * { chapter }
	{ \color{sealbrown} \vbox { \vskip 8.2pt \hbox to \l__zhvt_shukou_chapter_title_height_dim
		{ \hfil \fontsize{18pt}{18pt}\selectfont \l__zhvt_title_str \quad \l__zhvt_chapter_title_tl \hfil } } }
	\setdynamiccontents * { page~number } { \color{sealbrown} \quad \zhnum { page } }
}
\__zhvt_set_shukou_content:
\ExplSyntaxOff

\maketitle

% 单篇，无多卷目录，不渲染目录页

\juan{樗里子甘茂列传}{卷七十一~樗里子甘茂列传第十一}
\hspace*{2em}樗里子者。名疾。秦惠王之弟也。与惠王异母。母。韩女也。樗里子滑稽多智。秦人号曰﹁智囊﹂。\\
\hspace*{2em}秦惠王八年。爵樗里子右更。使将而伐曲沃。尽出其人。取其城。地入秦。秦惠王二十五年。使樗里子为将伐赵。虏赵将军庄豹。拔蔺。明年。助魏章攻楚。败楚将屈丐。取汉中地。秦封樗里子。号为严君。\\
\hspace*{2em}秦惠王卒。太子武王立。逐张仪。魏章。而以樗里子。甘茂为左右丞相。秦使甘茂攻韩。拔宜阳。使樗里子以车百乘入周。周以卒迎之。意甚敬。楚王怒。让周。以其重秦客。游腾为周说楚王曰。﹁知伯之伐仇犹。遗之广车。因随之以兵。仇犹遂亡。何则。无备故也。齐桓公伐蔡。号曰诛楚。其实袭蔡。今秦。虎狼之国。使樗里子以车百乘入周。周以仇犹。蔡观焉。故使长戟居前。彊弩在後。名曰卫疾。而实囚之。且夫周岂能无忧其社稷哉。恐一旦亡国以忧大王。﹂楚王乃悦。\\
\hspace*{2em}秦武王卒。昭王立。樗里子又益尊重。\\
\hspace*{2em}昭王元年。樗里子将伐蒲。蒲守恐。请胡衍。胡衍为蒲谓樗里子曰。﹁公之攻蒲。为秦乎。为魏乎。为魏则善矣。为秦则不为赖矣。夫卫之所以为卫者。以蒲也。今伐蒲入於魏。卫必折而从之。魏亡西河之外而无以取者。兵弱也。今并卫於魏。魏必彊。魏彊之日。西河之外必危矣。且秦王将观公之事。害秦而利魏。王必罪公。﹂樗里子曰。﹁柰何。﹂胡衍曰。﹁公释蒲勿攻。臣试为公入言之。以德卫君。﹂樗里子曰。﹁善。﹂胡衍入蒲。谓其守曰。﹁樗里子知蒲之病矣。其言曰必拔蒲。衍能令释蒲勿攻。﹂蒲守恐。因再拜曰。﹁愿以请。﹂因效金三百斤。曰。﹁秦兵苟退。请必言子於卫君。使子为南面。﹂故胡衍受金於蒲以自贵於卫。於是遂解蒲而去。还击皮氏。皮氏未降。又去。\\
\hspace*{2em}昭王七年。樗里子卒。葬于渭南章台之东。曰。﹁後百岁。是当有天子之宫夹我墓。﹂樗里子疾室在於昭王庙西渭南阴乡樗里。故俗谓之樗里子。至汉兴。长乐宫在其东。未央宫在其西。武库正直其墓。秦人谚曰。﹁力则任鄙。智则樗里。﹂\\
\hspace*{2em}甘茂者。下蔡人也。事下蔡史举先生。学百家之术。因张仪。樗里子而求见秦惠王。王见而说之。使将。而佐魏章略定汉中地。\\
\hspace*{2em}惠王卒。武王立。张仪。魏章去。东之魏。蜀侯煇。相壮反。秦使甘茂定蜀。还。而以甘茂为左丞相。以樗里子为右丞相。\\
\hspace*{2em}秦武王三年。谓甘茂曰。﹁寡人欲容车通三川。以窥周室。而寡人死不朽矣。﹂甘茂曰。﹁请之魏。约以伐韩。而令向寿辅行。﹂甘茂至。谓向寿曰。﹁子归。言之於王曰﹃魏听臣矣。然愿王勿伐﹄。事成。尽以为子功。﹂向寿归。以告王。王迎甘茂於息壤。甘茂至。王问其故。对曰。﹁宜阳。大县也。上党。南阳积之久矣。名曰县。其实郡也。今王倍数险。行千里攻之。难。昔曾参之处费。鲁人有与曾参同姓名者杀人。人告其母曰﹃曾参杀人﹄。其母织自若也。顷之。一人又告之曰﹃曾参杀人﹄。其母尚织自若也。顷又一人告之曰﹃曾参杀人﹄。其母投杼下机。逾墙而走。夫以曾参之贤与其母信之也。三人疑之。其母惧焉。今臣之贤不若曾参。王之信臣又不如曾参之母信著参也。疑臣者非特三人。臣恐大王之投杼也。始张仪西并巴蜀之地。北开西河之外。南取上庸。天下不以多张子而以贤先王。魏文侯令乐羊将而攻中山。三年而拔之。乐羊返而论功。文侯示之谤书一箧。乐羊再拜稽首曰。﹃此非臣之功也。主君之力也。﹄今臣。羁旅之臣也。樗里子。公孙奭二人者挟韩而议之。王必听之。是王欺魏王而臣受公仲侈之怨也。﹂王曰。﹁寡人不听也。请与子盟。﹂卒使丞相甘茂将兵伐宜阳。五月而不拔。樗里子。公孙奭果争之。武王召甘茂。欲罢兵。甘茂曰。﹁息壤在彼。﹂王曰。﹁有之。﹂因大悉起兵。使甘茂击之。斩首六万。遂拔宜阳。韩襄王使公仲侈入谢。与秦平。\\
\hspace*{2em}武王竟至周。而卒於周。其弟立。为昭王。王母宣太后。楚女也。楚怀王怨前秦败楚於丹阳而韩不救。乃以兵围韩雍氏。韩使公仲侈告急於秦。秦昭王新立。太后楚人。不肯救。公仲因甘茂。茂为韩言於秦昭王曰。﹁公仲方有得秦救。故敢扞楚也。今雍氏围。秦师不下殽。公仲且仰首而不朝。公叔且以国南合於楚。楚。韩为一。魏氏不敢不听。然则伐秦之形成矣。不识坐而待伐孰与伐人之利。﹂秦王曰。﹁善。﹂乃下师於殽以救韩。楚兵去。\\
\hspace*{2em}秦使向寿平宜阳。而使樗里子。甘茂伐魏皮氏。向寿者。宣太后外族也。而与昭王少相长。故任用。向寿如楚。楚闻秦之贵向寿。而厚事向寿。向寿为秦守宜阳。将以伐韩。韩公仲使苏代谓向寿曰。﹁禽困覆车。公破韩。辱公仲。公仲收国复事秦。自以为必可以封。今公与楚解地。封小令尹以杜阳。秦楚合。复攻韩。韩必亡。韩亡。公仲且躬率其私徒以阏於秦。愿公孰虑之也。﹂向寿曰。﹁吾合秦楚非以当韩也。子为寿谒之公仲。曰秦韩之交可合也。﹂苏代对曰。﹁愿有谒於公。人曰贵其所以贵者贵。王之爱习公也。不如公孙奭。其智能公也。不如甘茂。今二人者皆不得亲於秦事。而公独与王主断於国者何。彼有以失之也。公孙奭党於韩。而甘茂党於魏。故王不信也。今秦楚争彊而公党於楚。是与公孙奭。甘茂同道也。公何以异之。人皆言楚之善变也。而公必亡之。是自为责也。公不如与王谋其变也。善韩以备楚。如此则无患矣。韩氏必先以国从公孙奭而後委国於甘茂。韩。公之雠也。今公言善韩以备楚。是外举不僻雠也。﹂向寿曰。﹁然。吾甚欲韩合。﹂对曰。﹁甘茂许公仲以武遂。反宜阳之民。今公徒收之。甚难。﹂向寿曰。﹁然则奈何。武遂终不可得也。﹂对曰。﹁公奚不以秦为韩求颍川於楚。此韩之寄地也。公求而得之。是令行於楚而以其地德韩也。公求而不得。是韩楚之怨不解而交走秦也。秦楚争彊。而公徐过楚以收韩。此利於秦。﹂向寿曰。﹁柰何。﹂对曰。﹁此善事也。甘茂欲以魏取齐。公孙奭欲以韩取齐。今公取宜阳以为功。收楚韩以安之。而诛齐魏之罪。是以公孙奭。甘茂无事也。﹂\\
\hspace*{2em}甘茂竟言秦昭王。以武遂复归之韩。向寿。公孙奭争之。不能得。向寿。公孙奭由此怨。谗甘茂。茂惧。辍伐魏蒲阪。亡去。樗里子与魏讲。罢兵。\\
\hspace*{2em}甘茂之亡秦奔齐。逢苏代。代为齐使於秦。甘茂曰。﹁臣得罪於秦。惧而遯逃。无所容迹。臣闻贫人女与富人女会绩。贫人女曰。﹃我无以买烛。而子之烛光幸有馀。子可分我馀光。无损子明而得一斯便焉。﹄今臣困而君方使秦而当路矣。茂之妻子在焉。愿君以馀光振之。﹂苏代许诺。遂致使於秦。已。因说秦王曰。﹁甘茂。非常士也。其居於秦。累世重矣。自殽塞及至鬼谷。其地形险易皆明知之。彼以齐约韩魏反以图秦。非秦之利也。﹂秦王曰。﹁然则柰何。﹂苏代曰。﹁王不若重其贽。厚其禄以迎之。使彼来则置之鬼谷。终身勿出。﹂秦王曰。﹁善。﹂即赐之上卿。以相印迎之於齐。甘茂不往。苏代谓齐湣王曰。﹁夫甘茂。贤人也。今秦赐之上卿。以相印迎之。甘茂德王之赐。好为王臣。故辞而不往。今王何以礼之。﹂齐王曰。﹁善。﹂即位之上卿而处之。秦因复甘茂之家以市於齐。\\
\hspace*{2em}齐使甘茂於楚。楚怀王新与秦合婚而驩。而秦闻甘茂在楚。使人谓楚王曰。﹁愿送甘茂於秦。﹂楚王问於范蜎曰。﹁寡人欲置相於秦。孰可。﹂对曰。﹁臣不足以识之。﹂楚王曰。﹁寡人欲相甘茂。可乎。﹂对曰。﹁不可。夫史举。下蔡之监门也。大不为事君。小不为家室。以苟贱不廉闻於世。甘茂事之顺焉。故惠王之明。武王之察。张仪之辩。而甘茂事之。取十官而无罪。茂诚贤者也。然不可相於秦。夫秦之有贤相。非楚国之利也。且王前尝用召滑於越。而内行章义之难。越国乱。故楚南塞厉门而郡江东。计王之功所以能如此者。越国乱而楚治也。今王知用诸越而忘用诸秦。臣以王为钜过矣。然则王若欲置相於秦。则莫若向寿者可。夫向寿之於秦王。亲也。少与之同衣。长与之同车。以听事。王必相向寿於秦。则楚国之利也。﹂於是使使请秦相向寿於秦。秦卒相向寿。而甘茂竟不得复入秦。卒於魏。\\
\hspace*{2em}甘茂有孙曰甘罗。\\
\hspace*{2em}甘罗者。甘茂孙也。茂既死後。甘罗年十二。事秦相文信侯吕不韦。\\
\hspace*{2em}秦始皇帝使刚成君蔡泽於燕。三年而燕王喜使太子丹入质於秦。秦使张唐往相燕。欲与燕共伐赵以广河间之地。张唐谓文信侯曰。﹁臣尝为秦昭王伐赵。赵怨臣。曰。﹃得唐者与百里之地。﹄今之燕必经赵。臣不可以行。﹂文信侯不快。未有以彊也。甘罗曰。﹁君侯何不快之甚也。﹂文信侯曰。﹁吾令刚成君蔡泽事燕三年。燕太子丹已入质矣。吾自请张卿相燕而不肯行。﹂甘罗曰。﹁臣请行之。﹂文信侯叱曰。﹁去。我身自请之而不肯。女焉能行之。﹂甘罗曰。﹁大项橐生七岁为孔子师。今臣生十二岁於兹矣。君其试臣。何遽叱乎。﹂於是甘罗见张卿曰。﹁卿之功孰与武安君。﹂卿曰。﹁武安君南挫彊楚。北威燕。赵。战胜攻取。破城堕邑。不知其数。臣之功不如也。﹂甘罗曰。﹁应侯之用於秦也。孰与文信侯专。﹂张卿曰。﹁应侯不如文信侯专。﹂甘罗曰。﹁卿明知其不如文信侯专与。﹂曰。﹁知之。﹂甘罗曰。﹁应侯欲攻赵。武安君难之。去咸阳七里而立死於杜邮。今文信侯自请卿相燕而不肯行。臣不知卿所死处矣。﹂张唐曰。﹁请因孺子行。﹂令装治行。\\
\hspace*{2em}行有日。甘罗谓文信侯曰。﹁借臣车五乘。请为张唐先报赵。﹂文信侯乃入言之於始皇曰。﹁昔甘茂之孙甘罗。年少耳。然名家之子孙。诸侯皆闻之。今者张唐欲称疾不肯行。甘罗说而行之。今愿先报赵。请许遣之。﹂始皇召见。使甘罗於赵。赵襄王郊迎甘罗。甘罗说赵王曰。﹁王闻燕太子丹入质秦欤。﹂曰。﹁闻之。﹂曰。﹁闻张唐相燕欤。﹂曰。﹁闻之。﹂﹁燕太子丹入秦者。燕不欺秦也。张唐相燕者。秦不欺燕也。燕。秦不相欺者。伐赵。危矣。燕。秦不相欺无异故。欲攻赵而广河间。王不如赍臣五城以广河间。请归燕太子。与彊赵攻弱燕。﹂赵王立自割五城以广河间。秦归燕太子。赵攻燕。得上谷三十城。令秦有十一。\\
\hspace*{2em}甘罗还报秦。乃封甘罗以为上卿。复以始甘茂田宅赐之。\\
\hspace*{2em}太史公曰。樗里子以骨肉重。固其理。而秦人称其智。故颇采焉。甘茂起下蔡闾阎。显名诸侯。重彊齐楚。甘罗年少。然出一奇计。声称後世。虽非笃行之君子。然亦战国之策士也。方秦之彊时。天下尤趋谋诈哉。\\

% ===== 索隐述赞（司马贞《史记索隐》）· 统一样式 =====
% 题头「索隐述赞曰」五字并入反白框；正文每四句一列，定宽盒 \hfil 居列中；
% \fontsize{24}{41} 列距对齐全书界行（121+41k）；\clearpage 使述赞独占一页。
\clearpage
\shuzanhead{索隐述赞曰}\linebreak
{\shuzanbody\fontsize{24}{41}\selectfont
\hbox to \hsize{\hfil 严君名疾。厥号﹁智囊﹂。既亲且重。称兵外攘。\hfil}\linebreak
\hbox to \hsize{\hfil 甘茂并相。初佐魏章。始推向寿。乃攻宜阳。\hfil}\linebreak
\hbox to \hsize{\hfil 甘罗妙岁。卒起张唐。\hfil}%
\par}
\end{document}
